% =============================================================================
% ANHANG B: VOLLSTÄNDIGES FORMALES MODELL DES ZEITLICHEN KOORDINATIONSPARADOXONS
% =============================================================================

\documentclass[12pt]{article}
\usepackage[a4paper,left=2cm,right=2cm,top=2.5cm,bottom=2.5cm]{geometry}

% Schriftkonfiguration für XeLaTeX (wie in den vorherigen deutschen Anhängen)
\usepackage{fontspec}
\usepackage{polyglossia}
\setmainlanguage{german}

\usepackage{amsmath, amssymb, amsthm}
\usepackage{hyperref}
\usepackage{titlesec}
\usepackage{graphicx}
\usepackage{tikz}
\usepackage{caption}
\usepackage{booktabs}
\usepackage{array}

% TikZ-Bibliotheken
\usetikzlibrary{positioning, arrows.meta, shapes.geometric, calc}

\titleformat{\section}{\large\bfseries}{\thesection}{1em}{}
\titleformat{\subsection}{\normalsize\bfseries}{\thesubsection}{1em}{}
\titleformat{\subsubsection}{\normalsize\bfseries}{\thesubsubsection}{1em}{}

\renewcommand{\baselinestretch}{1.1}

\newtheorem{definition}{Definition}
\newtheorem{theorem}{Theorem}
\newtheorem{lemma}{Lemma}
\newtheorem{corollary}{Korollar}
\newtheorem{axiom}{Axiom}

\begin{document}
	\begin{titlepage}
		\begin{center}
			\vspace*{0cm}
			\Huge\textbf{Anhang B. YUCT\\
				ZEITLICHES KOORDINATIONSPARADOXON}
			
			\vspace{0cm}
			\LARGE\textit{Von informationstheoretischen Grundlagen zur experimentellen Verifikation \\
				Mathematische Auflösung des $K_{\text{eff}} > 1$-Paradoxons}
			
			\vspace{1cm}
			\Large
			Alexey V. Yakushev\\
			Yakushev Research\\
			YUCT Theorie: \url{https://yuct.org/}\\
			YPSDC Protokoll: \url{https://ypsdc.com/}\\
			
			
			\vspace{2cm}
			\textbf DOI: \url{https://doi.org/10.5281/zenodo.18444598}\\
			
			\vspace{1cm}
			
			\vspace{0cm}
			\large 
			Eine Kurzfassung dieser Arbeit wurde zuvor veröffentlicht und ist verfügbar unter\\ \url{https://doi.org/10.5281/zenodo.18362308}\\
			\vspace{1cm}
			März 2026 \\ \large V.2.1.
			
			\vspace{5cm}
			\textcopyright~2026 Yakushev Research. Alle Rechte vorbehalten.
			
			\newpage
			\vfill
			\normalsize
			\begin{minipage}{0.8\textwidth}
				\centering
				\textbf{Zusammenfassung:} \\ Dieses Dokument präsentiert das vollständige formale mathematische Modell des zeitlichen Koordinationsparadoxons in der Yakushev Unified Coordination Theory (YUCT). Das Paradoxon tritt auf, wenn die Koordinationseffizienz $K_{\text{eff}}$ den Wert 1 überschreitet und scheinbar die klassischen Grenzen der Informationsübertragung verletzt. Wir liefern rigorose Definitionen, Theoreme und Beweise, die zeigen, wie Vorabwissen durch vorherige Wörterbücher eine Koordination ermöglicht, die zeitlich paradox erscheint, aber mit den physikalischen Gesetzen konsistent bleibt. Zu den wichtigsten Ergebnissen gehören: (1) Formales Separations-Theorem für Kanalkapazität vs. Koordinationskapazität, (2) Mathematische Herleitung von $K_{\text{eff}} = R \cdot \eta$ mit $R = n/m$ als Wissenskompressionsverhältnis, (3) Exakte Bedingungen, unter denen $K_{\text{eff}}$ den Wert 1 überschreiten kann, ohne die Kausalität zu verletzen, (4) Mehrknoten-Koordinationstheoreme für skalierbare Systeme, (5) Entropische Interpretation der Koordinationseffizienz und (6) Experimentelle Protokolle zur Messung von $K_{\text{eff}}$ in realen Systemen.
			\end{minipage}
			
			\vspace{0.5cm}
			\normalsize
			\textbf{Schlüsselwörter:} Zeitliches Koordinationsparadoxon, YUCT formales Modell, $K_{\text{eff}} > 1$, Vorabwissen, vorherige Wörterbücher, informationstheoretische Koordination, Kanalkapazitätstrennung, YPSDC-Protokoll, Kausalitätserhaltung, Koordinationsentropie.
			
			\vspace{0cm}
			\normalsize
			
			\textbf{Anwendungsbereich:} Formale mathematische Grundlagen der zeitlichen Koordination \\
			\textbf{Wichtige Theoreme:} 8 formale Theoreme mit vollständigen Beweisen \\
			\textbf{Experimentelle Protokolle:} 3 Messmethoden für $K_{\text{eff}}$ \\
			
			\vspace{0.5cm}
			\normalsize
			
		\end{center}
	\end{titlepage}	
	\begin{center}
		\vspace*{0.5cm}
		\Huge\textbf{Zeitliches Koordinationsparadoxon}
		\vspace{1cm}
	\end{center}
	
	\section{Vollständiges formales Modell des zeitlichen Koordinationsparadoxons}
	\label{app:formal-model}
	
	\subsection{Formale Definitionen}
	\label{app:definitions}
	
	\subsubsection{Grundlegende Objekte}
	
	Ein Koordinationssystem sei definiert als ein Tupel $\mathcal{S} = (\mathcal{A}, \mathcal{O}, \mathcal{D}, \mathcal{C}, \mathcal{T})$, wobei:
	
	\begin{itemize}
		\item $\mathcal{A} = \{a_1, a_2, \dots, a_N\}$ eine endliche Menge möglicher Aktionen (Wissensaktivierungen) ist,
		\item $\mathcal{O} = \{O_1, O_2, \dots, O_M\}$ eine Menge von Beobachtern (Knoten) ist,
		\item $\mathcal{D} = \{\kappa_i \to a_i\}_{i=1}^N$ ein vorheriges Wörterbuch (bijektive Abbildung von Codes auf Aktionen) ist,
		\item $\mathcal{C}$ ein physikalischer Übertragungskanal mit spezifizierten Parametern ist,
		\item $\mathcal{T}$ eine Menge zeitlicher Nebenbedingungen ist.
	\end{itemize}
	
	\subsubsection{Informationstheoretische Maße}
	
	\begin{definition}[Informationsgehalt einer Aktion]
		Für jede Aktion $a \in \mathcal{A}$ erfordert ihre vollständige Beschreibung:
		\[
		I(a) = n \ \text{Bits},
		\]
		wobei $n$ die minimale Anzahl von Bits ist, die benötigt wird, um $a$ ohne Vorabwissen eindeutig zu beschreiben.
	\end{definition}
	
	\begin{definition}[Codelänge]
		Für jeden Code $\kappa \in \mathcal{K}$ (der Menge der Codes in $\mathcal{D}$):
		\[
		|\kappa| = m \ \text{Bits}, \quad \text{mit} \ m \ll n.
		\]
	\end{definition}
	
	\begin{definition}[Wissenskompressionsverhältnis]
		Das Wissenskompressionsverhältnis ist definiert als:
		\[
		R = \frac{n}{m} \gg 1.
		\]
	\end{definition}
	
	\subsection{Physikalisches Kanalmodell}
	\label{app:channel-model}
	
	\subsubsection{Fundamentale Nebenbedingungen}
	
	\begin{axiom}[Lichtgeschwindigkeitsgrenze]
		Für zwei beliebige Knoten $O_i, O_j \in \mathcal{O}$, die durch die Distanz $L_{ij}$ getrennt sind, erfüllt die Signalausbreitungsgeschwindigkeit:
		\[
		v_{\text{Signal}} \leq c,
		\]
		wobei $c$ die Lichtgeschwindigkeit im Vakuum ist.
	\end{axiom}
	
	\begin{axiom}[Kanalkapazität]
		Der Kanal $\mathcal{C}$ hat eine maximale Kanalkapazität, gegeben durch:
		\[
		C_{\text{max}} = \lim_{T \to \infty} \frac{1}{T} \max_{X} I(X; Y) \quad \text{[Bits/s]},
		\]
		wobei $X$ das Eingangssignal und $Y$ das Ausgangssignal ist.
	\end{axiom}
	
	\subsubsection{Modell der Übertragungszeit}
	
	Die Zeit, die benötigt wird, um eine Nachricht von $I$ Bits zu übertragen, ist:
	\[
	T_{\text{übertragen}}(I) = \frac{I}{C_{\text{eff}}} + \frac{L}{c} + \tau_{\text{proc}},
	\]
	wobei $C_{\text{eff}} \leq C_{\text{max}}$ die effektive Kanalkapazität und $\tau_{\text{proc}}$ die Verarbeitungsverzögerung ist.
	
	\subsection{Koordinationsprozesse}
	\label{app:coordination-processes}
	
	\subsubsection{Naive Koordination (ohne Wörterbuch)}
	
	\textbf{Szenario 1:} Knoten $O_1$ möchte, dass Knoten $O_2$ die Aktion $a$ ausführt.
	
	\begin{enumerate}
		\item $O_1$ formuliert die vollständige Beschreibung von $a$: $I(a) = n$ Bits.
		\item $O_1$ überträgt diese Beschreibung an $O_2$: Zeit $T_1 = T_{\text{übertragen}}(n)$.
		\item $O_2$ decodiert und versteht die Aktion: Zeit $T_2 = \alpha n$, wobei $\alpha > 0$.
		\item $O_2$ führt die Aktion aus: Zeit $T_3$.
		\item $O_2$ sendet eine Bestätigung: Zeit $T_4 = T_{\text{übertragen}}(p)$, wobei $p$ die Größe der Bestätigung ist.
	\end{enumerate}
	
	Gesamte naive Koordinationszeit:
	\[
	T_{\text{naiv}}(a) = T_1 + T_2 + T_3 + T_4.
	\]
	
	\begin{theorem}[Untere Schranke für naive Koordination]
		\label{thm:naive-lower-bound}
		Für jedes Paar von Knoten $O_1, O_2$ und jede Aktion $a$ ist die naive Koordinationszeit beschränkt durch:
		\[
		T_{\text{naiv}}(a) \geq \frac{n + p}{C_{\text{eff}}} + \frac{2L}{c} + \alpha n.
		\]
	\end{theorem}
	
	\subsubsection{Wörterbuchbasierte Koordination (YPSDC)}
	
	\textbf{Szenario 2:} Die Knoten $O_1$ und $O_2$ teilen ein vorheriges Wörterbuch $\mathcal{D}$.
	
	\begin{enumerate}
		\item $O_1$ wählt den Code $\kappa$, der der Aktion $a$ entspricht: $|\kappa| = m$ Bits.
		\item $O_1$ überträgt $\kappa$ an $O_2$: Zeit $T_1' = T_{\text{übertragen}}(m)$.
		\item $O_2$ schlägt die Aktion in $\mathcal{D}$ nach: Zeit $T_2' = \beta m$, mit $\beta \ll \alpha$.
		\item $O_2$ führt die Aktion aus: Zeit $T_3' = T_3$.
	\end{enumerate}
	
	Die kritische Beobachtung: $O_1$ kann so handeln, \emph{als ob} $O_2$ die Aktion $a$ bereits unmittelbar nach dem Senden von $\kappa$ ausgeführt hätte.
	
	Die effektive Koordinationszeit mit Wörterbuch ist:
	\[
	T_{\text{coord}}(a) = T_1' + T_2'.
	\]
	
	\begin{figure}[htbp]
		\centering
		\begin{tikzpicture}[node distance=2cm, auto, >=Stealth]
			\node (O1) [draw, rectangle, rounded corners, minimum width=2cm, minimum height=1cm] {Beobachter 1};
			\node (O2) [draw, rectangle, rounded corners, minimum width=2cm, minimum height=1cm, right=of O1] {Beobachter 2};
			\node (dict) [draw, ellipse, below=2cm of $(O1)!0.5!(O2)$] {Vorheriges Wörterbuch $\mathcal{D}$};
			\draw[->, thick] (O1) -- node[above] {kurzer Code $\kappa$} (O2);
			\draw[->, dashed] (O1) -- node[left, near start] {teilen} (dict);
			\draw[->, dashed] (O2) -- node[right, near start] {teilen} (dict);
			\node[below=0.3cm of O1, align=center] {bei $t_0$:\\ weiß, dass $a$ ausgeführt wird};
			\node[below=0.3cm of O2, align=center] {bei $t_0 + L/c$:\\ führt $a$ aus};
		\end{tikzpicture}
		\caption{Grundlegendes YPSDC-Koordinationsschema. Vorabwissen entsteht bei $t_0$ durch das gemeinsame Wörterbuch.}
		\label{fig:basic}
	\end{figure}
	
	\subsection{Das zeitliche Koordinationsparadoxon}
	\label{app:paradox}
	
	\begin{definition}[Vorabwissen]
		Im Moment $t_0$, wenn der Code $\kappa$ gesendet wird, besitzt Knoten $O_1$ Vorabwissen über die zukünftige Aktion von $O_2$:
		\[
		K_{\text{vorab}}(O_1, O_2, a, t_0) = \text{Wahr}
		\]
		genau dann, wenn es ein gemeinsames Wörterbuch $\mathcal{D}$ gibt, so dass $\mathcal{D}(\kappa) = a$.
	\end{definition}
	
	\begin{lemma}[Existenz von Vorabwissen]
		\label{lem:advance-knowledge}
		Für beliebige $O_1, O_2, a$ mit einem gemeinsamen Wörterbuch $\mathcal{D}$ gilt:
		\[
		K_{\text{vorab}}(O_1, O_2, a, t_0) = \text{Wahr}
		\]
		zum Zeitpunkt $t_0$ des Sendens von $\kappa$.
	\end{lemma}
	
	\begin{proof}
		Durch die Konstruktion von $\mathcal{D}$ garantiert das Senden von $\kappa$, dass $O_2$ die Aktion $a$ zum Zeitpunkt $t_0 + T_{\text{coord}}(a)$ ausführen wird. Daher kann $O_1$ zum Zeitpunkt $t_0$ den zukünftigen Zustand von $O_2$ mit Sicherheit vorhersagen.
	\end{proof}
	
	\subsection{Informationskapazitäten der Koordination}
	\label{app:capacities}
	
	\begin{definition}[Informationskapazität des Kanals]
		\[
		C_{\text{Kanal}} = C_{\text{eff}} \quad \text{[Bits/s]}.
		\]
	\end{definition}
	
	\begin{definition}[Koordinationsinformationskapazität]
		\[
		C_{\text{coord}} = \lim_{T \to \infty} \frac{1}{T} \sum_{t=1}^{T} I(a_t) \quad \text{[Bits/s]},
		\]
		wobei $a_t$ die zum Zeitpunkt $t$ aktivierte Aktion ist.
	\end{definition}
	
	\begin{theorem}[Kapazitätstrennungs-Theorem]
		\label{thm:capacity-separation}
		Für ein Koordinationssystem $\mathcal{S}$ mit Wissenskompressionsverhältnis $R = n/m$ gilt:
		\[
		C_{\text{coord}} = R \cdot C_{\text{Kanal}} \cdot \eta,
		\]
		wobei $\eta = \frac{T_{\text{Kanal}}}{T_{\text{coord}}} \leq 1$ der Zeiteffizienzfaktor ist.
	\end{theorem}
	
	\begin{proof}
		Während eines Zeitintervalls $\Delta T$ kann der Kanal übertragen:
		\[
		N_{\text{Codes}} = \frac{C_{\text{Kanal}} \cdot \Delta T}{m} \quad \text{Codes}.
		\]
		Jeder Code aktiviert eine Aktion mit Informationsgehalt $n$ Bits, also beträgt die gesamte koordinierte Information:
		\[
		I_{\text{gesamt}} = N_{\text{Codes}} \cdot n = \frac{C_{\text{Kanal}} \cdot \Delta T}{m} \cdot n.
		\]
		Folglich:
		\[
		C_{\text{coord}} = \frac{I_{\text{gesamt}}}{\Delta T} = C_{\text{Kanal}} \cdot \frac{n}{m} = R \cdot C_{\text{Kanal}}.
		\]
		Die Berücksichtigung von Zeitverzögerungen (z. B. Verarbeitung, Ausbreitung) führt den Effizienzfaktor $\eta$ ein.
	\end{proof}
	
	\begin{figure}[htbp]
		\centering
		\begin{tikzpicture}[node distance=2.5cm, auto, >=Stealth]
			\node (O1) [draw, rectangle, rounded corners, minimum width=2.2cm, minimum height=1.2cm] {Beobachter 1 $O_1$};
			\node (O2) [draw, rectangle, rounded corners, minimum width=2.2cm, minimum height=1.2cm, right=of O1] {Beobachter 2 $O_2$};
			\node (dict) [draw, ellipse, below=2.5cm of $(O1)!0.5!(O2)$] {Vorheriges Wörterbuch $\mathcal{D} = \{\kappa_i \to \mathcal{A}_i\}$};
			\draw[->, thick] (O1) -- node[above] {$\kappa_1$} (O2);
			\draw[->, thick] (O2) -- node[below] {$\kappa_2$ (Bestätigung)} (O1);
			\draw[->, dashed] (O1) -- (dict);
			\draw[->, dashed] (O2) -- (dict);
			\node[below=0.2cm of O1, align=center] {Kennt $\mathcal{A}_2$ bei $t_0$};
			\node[below=0.2cm of O2, align=center] {Führt $\mathcal{A}_2$ bei $t_0+L/c$ aus};
		\end{tikzpicture}
		\caption{Duplex-Koordination mit vorherigem Wörterbuch. Beide Knoten teilen dasselbe Wörterbuch, was beidseitiges Vorabwissen ermöglicht.}
		\label{fig:duplex}
	\end{figure}
	
	\subsection{Der Koordinationseffizienzfaktor \(K_{\mathrm{eff}}\) und die Kapazitätstrennung}
	
	\begin{definition}[Koordinationseffizienz]
		Die \emph{Koordinationseffizienz} eines YPSDC-Systems ist
		\begin{equation}
			K_{\mathrm{eff}} = \frac{H(\mathcal{A})}{H(\mathcal{K})},
			\label{eq:Kdef}
		\end{equation}
		wobei \(H(\mathcal{A})\) die Entropie des Aktionsraums und \(H(\mathcal{K})\) die Entropie der Indexverteilung ist.
	\end{definition}
	
	\begin{definition}[Koordinationskapazität]
		Die \emph{Koordinationskapazität} \(C_{\mathrm{coord}}\) ist die maximale Rate, mit der der Empfänger sinnvolle Information aus dem Wörterbuch extrahieren kann:
		\begin{equation}
			C_{\mathrm{coord}} = \lim_{\Delta t \to \infty} \frac{1}{\Delta t} \sum_{t=1}^{N_{\mathrm{Aktionen}}} I(A_{i_t}),
		\end{equation}
		wobei \(A_{i_t}\) die während der $t$-ten Übertragung aktivierte Aktion ist.
	\end{definition}
	
	\begin{theorem}[Kapazitätstrennung]
		\label{thm:capacity-separation}
		Für ein YPSDC-System mit Koordinationseffizienz \(K_{\mathrm{eff}}\) und einem Overhead-Faktor \(\eta \le 1\) gilt
		\begin{equation}
			\boxed{C_{\mathrm{coord}} = K_{\mathrm{eff}} \cdot C_{\mathrm{Kanal}} \cdot \eta}.
			\label{eq:capacity-separation}
		\end{equation}
	\end{theorem}
	
	\begin{proof}
		Über ein Zeitintervall \(\Delta T\) kann der Kanal höchstens \(C_{\mathrm{Kanal}} \Delta T\) Bits übertragen. Bei festen Indexlängen von \(\ell\) Bits ist die Anzahl übertragener Indizes \(N = \lfloor C_{\mathrm{Kanal}} \Delta T / \ell \rfloor\). Jeder Index aktiviert eine Aktion mit durchschnittlich \(H(\mathcal{A})\) Bits Information. Die gesamte sinnvolle gelieferte Information ist \(I_{\mathrm{gesamt}} = N \cdot H(\mathcal{A}) \approx \frac{C_{\mathrm{Kanal}} \Delta T}{\ell} H(\mathcal{A})\). Die mittlere Rate ist \(C_{\mathrm{coord}} = \frac{I_{\mathrm{gesamt}}}{\Delta T} = C_{\mathrm{Kanal}} \cdot \frac{H(\mathcal{A})}{\ell} = C_{\mathrm{Kanal}} \cdot K_{\mathrm{eff}}\). Der Faktor \(\eta\) berücksichtigt Verarbeitungs- und Wörterbuchzugriffs-Overheads.
	\end{proof}
	
	\begin{theorem}[Umkehrung]
		Für jedes Schema, das ein festes Wörterbuch \(D\) der Größe \(M = 2^\ell\) verwendet und Indizes über einen DMC (diskreten gedächtnislosen Kanal) mit der Kapazität \(C_{\mathrm{Kanal}}\) überträgt, erfüllt die Koordinationskapazität
		\[
		C_{\mathrm{coord}} \le K_{\mathrm{eff}} \cdot C_{\mathrm{Kanal}}.
		\]
	\end{theorem}
	
	\begin{proof}
		Die Indexquelle habe die Entropie \(H(\mathcal{K})\). Nach der Datenverarbeitungsungleichung gilt \(I(\kappa;\hat{\kappa}) \le I(X;Y) \le C_{\mathrm{Kanal}}\). Für zuverlässigen Betrieb muss der Index mit verschwindendem Fehler decodiert werden, was eine Indexrate \(R_{\mathcal{K}} \le C_{\mathrm{Kanal}}\) erfordert (Shannons Rauschcodierungstheorem). Die mittlere sinnvolle Informationsrate ist \(R_{\mathcal{A}} = R_{\mathcal{K}} \cdot K_{\mathrm{eff}} \le K_{\mathrm{eff}} \cdot C_{\mathrm{Kanal}}\). Somit ist \(C_{\mathrm{coord}} \le K_{\mathrm{eff}} \cdot C_{\mathrm{Kanal}}\).
	\end{proof}
	
	\subsection{Entropische Interpretation der Koordination}
	\label{app:entropy}
	
	\begin{definition}[Wörterbuchentropie]
		Die Entropie des Wörterbuchs $\mathcal{D}$ ist:
		\[
		H(\mathcal{D}) = - \sum_{i=1}^{N} p_i \log_2 p_i,
		\]
		wobei $p_i$ die Wahrscheinlichkeit ist, die Aktion $a_i$ zu verwenden.
	\end{definition}
	
	\begin{lemma}[Maximale Koordinationskapazität]
		Die maximale Koordinationskapazität ist beschränkt durch:
		\[
		C_{\text{coord}}^{\text{max}} = H(\mathcal{D}) \cdot \nu_{\text{max}},
		\]
		wobei $\nu_{\text{max}} = 1 / T_{\text{coord}}^{\text{min}}$ die maximale Koordinationsfrequenz ist.
	\end{lemma}
	
	\begin{definition}[Koordinationsentropie]
		Die Änderung der Entropie des Systems durch Koordination ist:
		\[
		\Delta S_{\text{coord}} = k_B \ln(K_{\mathrm{eff}}),
		\]
		wobei $k_B$ die Boltzmann-Konstante ist.
	\end{definition}
	
	Interpretation: Ein Anstieg von $K_{\mathrm{eff}}$ entspricht einer Abnahme der Entropie (Zunahme der Ordnung) des koordinierten Systems.
	
	\subsection{Testbare Vorhersagen aus dem formalen Modell}
	\label{app:predictions}
	
	\begin{enumerate}
		\item \textbf{Quantisierung von $K_{\mathrm{eff}}$:} Für optimal entworfene Systeme gilt $K_{\mathrm{eff}} = 2^k$ für eine ganze Zahl $k \in \mathbb{N}$, wobei $k = \log_2 R$ unter idealen Bedingungen.
		
		\item \textbf{Skalierung mit der Systemgröße:} Für ein System mit $M$ Knoten gilt $K_{\mathrm{eff}}(M) \propto M \log_2 R$.
		
		\item \textbf{Fundamentale Grenze:} Es existiert eine fundamentale Grenze, gegeben durch:
		\[
		K_{\mathrm{eff}}^{\text{max}} = \frac{c}{\Delta x \cdot \Delta t},
		\]
		wobei $\Delta x$ die minimale räumliche Auflösung und $\Delta t$ die minimale zeitliche Auflösung des Systems ist.
	\end{enumerate}
	
	\subsection{Experimentelles Protokoll zur Messung von $K_{\mathrm{eff}}$}
	\label{app:experiment}
	
	\begin{enumerate}
		\item \textbf{Messen Sie $C_{\text{Kanal}}$:} Verwenden Sie Standard-Netzwerkmesswerkzeuge (z. B. iperf, ping), um die effektive Kanalkapazität zu schätzen.
		
		\item \textbf{Messen Sie $C_{\text{coord}}$:}
		\begin{itemize}
			\item Definieren Sie einen repräsentativen Satz von Aktionen $\mathcal{A}$.
			\item Messen Sie die Gesamtzeit $T_{\text{gesamt}}$, um eine Reihe koordinierter Aktionen auszuführen.
			\item Berechnen Sie $C_{\text{coord}} = \frac{\sum I(a_i)}{T_{\text{gesamt}}}$.
		\end{itemize}
		
		\item \textbf{Berechnen Sie $K_{\mathrm{eff}}$:}
		\[
		K_{\mathrm{eff}} = \frac{C_{\text{coord}}}{C_{\text{Kanal}}}.
		\]
	\end{enumerate}
	
	Erwartete Bereiche für verschiedene Systeme:
	\begin{itemize}
		\item Menschliche Befehlssysteme: $K_{\mathrm{eff}} \approx 10^2 - 10^3$.
		\item Computernetzwerke: $K_{\mathrm{eff}} \approx 10^3 - 10^6$.
		\item Biologische Systeme: $K_{\mathrm{eff}} \approx 10^6 - 10^9$.
	\end{itemize}
	
	\begin{table}[htbp]
		\centering
		\caption{Fundamentale Trennung: Datenübertragung vs. Koordination von Zuständen.}
		\label{tab:separation}
		\begin{tabular}{lp{6cm}p{6cm}}
			\toprule
			\textbf{Aspekt} & \textbf{Datenübertragung} & \textbf{Koordination} \\
			\midrule
			Physikalisches Signal & Ja & Nein (Zustandswissen) \\
			Geschwindigkeitsgrenze & $v \leq c$ & $v_{\text{coord}} \to \infty$* \\
			Information & $I > 0$ & $K_{\text{eff}} > I$ \\
			Energie & $E > 0$ & $\Delta S_{\text{coord}}$ \\
			\bottomrule
			\multicolumn{3}{l}{*im Sinne einer instantanen Zustandsangleichung durch Vorabwissen.}
		\end{tabular}
	\end{table}
	
	\subsection{Fazit des formalen Modells}
	
	Das in diesem Anhang präsentierte formale mathematische Modell etabliert streng:
	
	\begin{enumerate}
		\item Die Informationskapazität der Koordination $C_{\text{coord}}$ und die Kanalkapazität $C_{\text{Kanal}}$ sind verschiedene physikalische Größen.
		
		\item Das zeitliche Koordinationsparadoxon: A-priori-Wörterbücher erlauben es einem Knoten, Vorabwissen über zukünftige Aktionen anderer Knoten zu besitzen, bevor diese tatsächlich ausgeführt werden.
		
		\item Die quantitative Beziehung:
		\[
		K_{\mathrm{eff}} = \frac{C_{\text{coord}}}{C_{\text{Kanal}}} = R \cdot \eta \gg 1,
		\]
		wobei $R = n/m \gg 1$ das Wissenskompressionsverhältnis ist.
		
		\item Die fundamentale Einschränkung: Das Wachstum von $K_{\mathrm{eff}}$ wird nur durch die Komplexität der Erstellung und Aufrechterhaltung des Wörterbuchs beschränkt, nicht durch physikalische Gesetze der Informationsübertragung.
	\end{enumerate}
	
	Dieses Modell liefert eine strenge mathematische Grundlage zur Analyse und zum Entwurf hocheffizienter Koordinationssysteme in Physik, Biologie, Soziologie und Technologie.
	
	
	\section{Experimentelle Protokolle und Gedankenexperiment zum zeitlichen Koordinationsparadoxon}
	\label{app:experimental-protocols}
	
	Das formale Modell des zeitlichen Koordinationsparadoxons (Abschnitte \ref{app:formal-model}--\ref{app:keff}) liefert ein klares mathematisches Framework, aber seine praktische Verifikation erfordert konkrete experimentelle Protokolle. In diesem Abschnitt schlagen wir mehrere Experimente vor – von einem einfachen Gedankenexperiment bis zu realisierbaren Labortests – die die Koordinationseffizienz \(K_{\mathrm{eff}}\) direkt messen und das Paradoxon des Vorabwissens demonstrieren.
	
	\subsection{Gedankenexperiment: Die 2FA-Analogie neu betrachtet}
	\label{subsec:thought-2fa}
	
	Betrachten Sie zwei Beobachter, Alice und Bob, die ein kryptografisches Wörterbuch \( \mathcal{D} \) teilen (z. B. ein Einmal-Pad von \(M\) vorab geteilten Schlüsseln). Alice möchte, dass Bob eine von \(M\) möglichen Aktionen ausführt (z. B. eine Geldsumme überweisen, eine Rakete starten oder eine bestimmte Nachricht senden). In einem naiven Protokoll würde Alice die vollständige Beschreibung der Aktion senden, was \(n\) Bits erfordert. Mit dem gemeinsamen Wörterbuch sendet sie nur einen kurzen Index \(\kappa\) der Länge \(\ell = \log_2 M\) Bits.
	
	\textbf{Aufbau:}
	\begin{itemize}
		\item Alice und Bob sind durch eine bekannte Distanz \(L\) getrennt (z. B. 1 km).
		\item Der Kommunikationskanal hat eine gemessene Bandbreite und eine Ausbreitungsverzögerung \(L/c\).
		\item Das Wörterbuch \(\mathcal{D}\) wird offline über einen vertrauenswürdigen Boten oder über einen vorherigen sicheren Kanal verteilt (dies kann beliebig lange dauern, wird aber vor dem Experiment abgeschlossen).
		\item Ein dritter Beobachter, Charlie, ausgestattet mit zwei synchronisierten Atomuhren, zeichnet die genauen Zeiten auf, zu denen Alice die Übertragung initiiert und Bob die Aktion abschließt.
	\end{itemize}
	
	\textbf{Verfahren:}
	\begin{enumerate}
		\item Alice wählt eine Aktion \(a_i\) mit vollständiger Beschreibungslänge \(n\).
		\item Zum Zeitpunkt \(t_0\) sendet sie den entsprechenden Index \(\kappa_i\) an Bob.
		\item Der Index breitet sich mit Lichtgeschwindigkeit \(c\) aus und trifft bei Bob zum Zeitpunkt \(t_0 + L/c + \tau_{\text{tx}}\) ein, wobei \(\tau_{\text{tx}} = \ell / C_{\text{eff}}\).
		\item Bob schlägt die Aktion sofort im Wörterbuch nach (Nachschlagezeit \(\tau_{\text{Nachschlag}} \ll \tau_{\text{tx}}\)) und führt sie zum Zeitpunkt \(t_0 + L/c + \tau_{\text{tx}} + \tau_{\text{Nachschlag}}\) aus.
	\end{enumerate}
	
	\textbf{Beobachtung des Paradoxons:}
	Im Moment \(t_0\) \emph{weiß} Alice, dass Bob die Aktion \(a_i\) zu einem zukünftigen Zeitpunkt \(t_0 + L/c + \tau_{\text{tx}} + \tau_{\text{Nachschlag}}\) ausführen wird. Dieses Vorabwissen basiert nicht auf überlichtschneller Signalisierung, sondern auf dem vorab geteilten Wörterbuch. Charlie, der das Wörterbuch nicht besitzt, kann die Aktion erst nach der Übertragung des Index vorhersagen. Somit wird das Paradoxon aufgelöst, indem Information über das Wörterbuch (offline) vom Index (online) getrennt wird.
	
	\textbf{Quantitatives Maß:}
	Die Koordinationseffizienz ist
	\[
	K_{\mathrm{eff}} = \frac{n}{\ell} \cdot \frac{\tau_{\text{naiv}}}{\tau_{\text{tatsächlich}}},
	\]
	wobei \(\tau_{\text{naiv}}\) die Zeit ist, die zum Senden der vollständigen Beschreibung benötigt wird (einschließlich Ausbreitung und Verarbeitung). Im Grenzfall schneller Nachschlagezeiten und vernachlässigbarer Übertragungszeit für den Index gilt \(K_{\mathrm{eff}} \approx n/\ell\), was durch Vergrößerung des Wörterbuchs beliebig groß gemacht werden kann.
	
	\subsection{Realisierbares Laborexperiment: Verzögerte-Auswahl-Kommunikation}
	\label{subsec:lab-delayed-choice}
	
	\textbf{Ziel:} Direkte Messung von \(K_{\mathrm{eff}}\) in einem kontrollierten digitalen Kommunikationssystem und Verifikation, dass \(K_{\mathrm{eff}} > 1\) ohne Verletzung der Kausalität erreicht werden kann.
	
	\textbf{Apparatur:}
	\begin{itemize}
		\item Zwei Computer (Alice und Bob), verbunden durch eine Ethernet-Verbindung mit bekannter Kapazität und einer variablen Verzögerungsstrecke (z. B. eine lange Glasfaser), um eine Distanz zu simulieren.
		\item Ein vorab geteiltes Wörterbuch von \(M\) Nachrichten, jede mit einer Länge von \(n\) Bits (z. B. \(M = 2^{20}\), \(n = 10^6\) Bits).
		\item Hochpräzise Zeitstempel (Nanosekundenauflösung) unter Verwendung von GPS-synchronisierten Uhren.
		\item Software zur Messung der Zeit zwischen dem Senden eines Index und dem Empfang der Bestätigung der abgeschlossenen Aktion.
	\end{itemize}
	
	\textbf{Protokoll:}
	\begin{enumerate}
		\item Kalibrierungsphase: Messen Sie die Umlaufzeit für das Senden einer vollständigen Nachricht (naives Protokoll) über dieselbe Verbindung. Notieren Sie \(T_{\text{naiv}}\).
		\item Koordinationsphase: Für jeden Durchlauf wählt Alice zufällig eine Nachricht aus, sendet ihren Index und zeichnet die Zeit \(T_{\text{coord}}\) auf, bis Bob die Ausführung bestätigt.
		\item Wiederholen Sie dies für verschiedene Wörterbuchgrößen \(M\) und Nachrichtenlängen \(n\). Berechnen Sie \(K_{\mathrm{eff}} = T_{\text{naiv}} / T_{\text{coord}}\).
		\item Messen Sie auch die Kanalkapazität \(C_{\text{Kanal}}\) und die Wörterbuch-Nachschlagezeit, um das Kapazitätstrennungs-Theorem zu verifizieren.
	\end{enumerate}
	
	\textbf{Erwartetes Ergebnis:} Für ausreichend große Wörterbücher wird \(K_{\mathrm{eff}}\) deutlich größer als 1 sein, was zeigt, dass die effektive Rate sinnvoller Informationsübertragung die Kanalkapazität überschreiten kann. Das Experiment bestätigt auch, dass die Indexausbreitung \(v \le c\) gehorcht; die Kausalität bleibt erhalten, weil das Wörterbuch offline verteilt wurde.
	
	\subsection{Soziales Koordinationsexperiment: Menschliche Reaktionszeit mit gemeinsamem Kontext}
	\label{subsec:social-experiment}
	
	\textbf{Ziel:} Demonstration des zeitlichen Koordinationsparadoxons in einem menschlichen sozialen Umfeld, analog zum 2FA-Gedankenexperiment.
	
	\textbf{Aufbau:}
	\begin{itemize}
		\item Zwei Gruppen von Probanden (z. B. Studenten) werden auf eine Menge von \(M\) komplexen Befehlen trainiert (z. B. „ein Haus zeichnen“, „ein Rätsel lösen“, „ein Gedicht aufsagen“). Jeder Befehl benötigt mehrere Sekunden zur Ausführung, und seine vollständige Beschreibung (Text + Anweisungen) umfasst etwa \(n\) Bits.
		\item Die Trainingsphase etabliert ein gemeinsames „Wörterbuch“.
		\item Während des Experiments wird einer Gruppe (dem „Sender“) ein Befehl gezeigt, und sie muss ihn der anderen Gruppe (dem „Empfänger“) nur mit einem einzigen kurzen Wort (dem Index) übermitteln, das sie während des Trainings vereinbart haben.
		\item Die Distanz zwischen den Gruppen wird variiert (z. B. verschiedene Räume, Gebäude), um eine messbare Ausbreitungsverzögerung für Schall- oder Lichtsignale einzuführen.
		\item Die Zeit vom Moment, in dem der Sender den Befehl erhält, bis zum Moment, in dem der Empfänger die Aktion abschließt, wird gemessen.
	\end{itemize}
	
	\textbf{Kontrolle:} Dieselben Befehle werden auf naive Weise übermittelt (z. B. durch detaillierte Beschreibung der gesamten Aktion) ohne vorheriges Training.
	
	\textbf{Messungen:} Vergleichen Sie die Koordinationszeiten mit und ohne vorheriges Wörterbuch. Berechnen Sie \(K_{\mathrm{eff}}\) als Verhältnis der naiven Zeit zur koordinierten Zeit. Zeigen Sie, dass selbst wenn sich der Index mit endlicher Geschwindigkeit (Schall- oder Lichtgeschwindigkeit) ausbreitet, die effektive Koordination im Vergleich zur naiven Methode nahezu augenblicklich erscheinen kann.
	
	\textbf{Erwartetes Ergebnis:} \(K_{\mathrm{eff}}\) wird mit der Komplexität der Befehle und der Größe des gemeinsamen Wörterbuchs skalieren, was die theoretischen Vorhersagen bestätigt.
	
	\subsection{Experimentelle Verifikation des Kapazitätstrennungs-Theorems}
	\label{subsec:capacity-theorem}
	
	Das Kapazitätstrennungs-Theorem (Theorem \ref{thm:capacity-separation}) besagt, dass \(C_{\text{coord}} = K_{\mathrm{eff}} \cdot C_{\text{Kanal}} \cdot \eta\). Zur Verifikation müssen wir sowohl \(C_{\text{coord}}\) als auch \(C_{\text{Kanal}}\) unabhängig messen.
	
	\textbf{Methode:}
	\begin{itemize}
		\item Verwenden Sie denselben Aufbau wie in Abschnitt \ref{subsec:lab-delayed-choice}, variieren Sie jedoch die Kanalbandbreite (z. B. durch Einführung einer künstlichen Ratenbegrenzung).
		\item Messen Sie für jede Bandbreiteneinstellung die maximale Rate, mit der Indizes zuverlässig übertragen werden können (\(C_{\text{Kanal}}\)).
		\item Messen Sie gleichzeitig die Rate, mit der sinnvolle Aktionen abgeschlossen werden (\(C_{\text{coord}}\)), indem Sie die Anzahl der erfolgreich pro Sekunde aktivierten Aktionen zählen.
		\item Tragen Sie \(C_{\text{coord}}\) gegen \(C_{\text{Kanal}}\) für verschiedene Wörterbuchgrößen auf. Die Steigung sollte \(K_{\mathrm{eff}}\) sein, und der Achsenabschnitt sollte Null sein (bis auf den Zeiteffizienzfaktor \(\eta\)).
	\end{itemize}
	
	\subsection{Interpretation und Implikationen}
	
	Diese Experimente demonstrieren direkt die Kernidee von YPSDC: die Trennung von offline-Wörterbuchverteilung und online-Indexübertragung ermöglicht eine Koordinationseffizienz \(K_{\mathrm{eff}} > 1\), ohne die Kausalität zu verletzen. Das Paradoxon des Vorabwissens wird aufgelöst, weil das Wörterbuch vorab geteilt wird und der Index selbst keine neue Information trägt, die über das Verweisen auf einen vorhandenen Eintrag hinausgeht.
	
	Darüber hinaus liefern die Experimente quantitative Messungen von \(K_{\mathrm{eff}}\) unter kontrollierten Bedingungen, was die Kalibrierung der Theorie und den Vergleich mit realen Systemen (z. B. Internetprotokollen, Finanznetzwerken, biologischer Signalübertragung) ermöglicht. Sie dienen auch als Lehrmittel, um die kontraintuitive, aber kausale Natur hocheffizienter Koordination zu veranschaulichen.
	
	\subsection{Fazit des experimentellen Abschnitts}
	
	Durch die Implementierung dieser Protokolle – von einem einfachen Gedankenexperiment bis zu Labor- und Sozialtests – können Forscher das zeitliche Koordinationsparadoxon und das Kapazitätstrennungs-Theorem empirisch validieren. Die Ergebnisse werden bestätigen, dass \(K_{\mathrm{eff}}\) nicht nur ein theoretisches Konstrukt ist, sondern eine messbare Größe, die Eins überschreiten kann, was den Weg für praktische Anwendungen in Kommunikation, verteiltem Rechnen und kollektiver Entscheidungsfindung ebnet.
	
\end{document}